%% ======================================================================
%%  Chapter I -- Principal ideal rings
%% ======================================================================
\clearpage
\folio{13}%
\chapterhead{1}{I}{Principal ideal rings}{PRINCIPAL IDEAL RINGS}

%% ----------------------------------------------------------------------
\sechead{1}{Divisibility in principal ideal rings}%
        {DIVISIBILITY IN PRINCIPAL IDEAL RINGS}

Let $A$ be an integral domain, $K$ its field of fractions, $x$ and $y$
elements of $K$. We shall say that $x$ \emph{divides} $y$ if there exists
$a\in A$ such that $y=ax$. Equivalently, we say $x$ is a \emph{divisor} of
$y$, $y$ is a \emph{multiple} of $x$; notation $x\mid y$. This relation
between the elements of $K$ depends in an essential manner on the ring $A$;
if there is any confusion possible, we speak of divisibility in $K$ with
respect to $A$.

Given $x\in K$ we write $Ax$ for the set of multiples of $x$. Thus we may
write $y\in Ax$ in place of $x\mid y$, or $Ay\subset Ax$. The set $Ax$ is
called a \emph{principal fractional ideal} of $K$ with respect to $A$; if
$x\in A$, $Ax$ is the (ordinary) principal ideal of $A$ generated by $x$. As
the relation of divisibility, $x\mid y$, is equivalent to the order relation
$Ay\subset Ax$, divisibility possesses the following two properties
associated with order relations.
\[
x\mid x;\quad\text{if }x\mid y\text{ and }y\mid z,\text{ then }x\mid z.
\]

On the other hand, if $x\mid y$ and $y\mid x$, one cannot in general conclude
that $x=y$; one has only that $Ax=Ay$, which (if $y\neq0$) means that the
quotient $xy^{-1}$ is an invertible element of $A$; in this case $x$ and $y$
are called \emph{associates}; they are indistinguishable from the point of
view of divisibility.

\begin{statement}{Example}
The elements of $K$ which are associates of $1$ are the elements invertible
in $A$; they are often called the \emph{units} in $A$; they form a group
under multiplication, and we shall denote this group $A^{*}$. The
determination of the units in a ring $A$ is an interesting problem which we
shall treat in the case where $A$ is the ring of integers in a number field
(see \hyperref[chap:IV]{Chapter IV}). Here are some simple examples:
\begin{alist}
\item If $A$ is a field, $A^{*}=A-(0)$.
\item If $A=\mathbf{Z}$, $A^{*}=\{+1,-1\}$.
\item The units in the ring of polynomials $B=A[X_{1},\ldots,X_{n}]$, $A$ an
integral domain, are the invertible constants; in other words $B^{*}=A^{*}$.
\item The units in the ring of formal power series $A[[X_{1},\ldots,X_{n}]]$
are the formal power series whose constant term is invertible.
\end{alist}
\end{statement}

\begin{definition}\label{def:I.1.1}
A ring $A$ is called a \emph{principal ideal ring} if it is an integral
domain and if every ideal of $A$ is principal.
\end{definition}

We know that the ring $\mathbf{Z}$ of rational integers is a principal ideal
ring. (Recall that any ideal $\mathfrak{a}\neq(0)$ of $\mathbf{Z}$ contains a
least integer $b>0$. Dividing $x\in\mathfrak{a}$ by $b$ and using the fact
that $\mathbf{Z}$ is Euclidean, one sees that $x$ is a multiple of $b$.) If
$K$ is a field we know that the ring $K[X]$ of polynomials in one variable
over $K$ is a principal ideal ring (same method of proof: take a non-zero
polynomial $b(X)$ of lowest degree in the given ideal
$\mathfrak{a}\neq(0)$ and make use of the fact that $K[X]$ is a Euclidean
ring, i.e.\ the remainder under division of an arbitrary element of
$\mathfrak{a}$ by $b(X)$ must be of lower degree than $b(X)$ or zero, which
implies zero). This general method shows that any ``Euclidean ring'' (see
\hyperref[bib:1]{[1]}, Chapter VIII, §1, exercises, or \hyperref[bib:9]{[9]} Chapter 1) is a principal ideal
ring. If $K$ is a field, it is easy to see that any non-zero ideal in the
ring of formal power series $A=K[[X]]$ is of the form $AX^{n}$ with
$n\geqslant0$, so that $A=K[[X]]$ is a principal ideal ring too.

Divisibility in the field of fractions $K$ of a principal ideal ring $A$ is
particularly simple. We shall review it briefly.

\medskip
\noindent I.\quad Two arbitrary elements $u,v$ of $K$ have a \emph{greatest
common divisor} (gcd), i.e.\ an element $d$ for which the relations
\begin{equation}
\text{``}x\mid u\text{ and }x\mid v\text{''}\quad\text{and}\quad
\text{``}x\mid d\text{''}
\label{eq:I.1.1}
\end{equation}
are equivalent. This means the same thing as the assertion that $Au$ and
$Av$ have a least upper bound in the partially ordered set of fractional
ideals. This least upper bound is $Au+Av$, which is itself a principal
fractional ideal, since the ring $A$ is a principal ideal ring (all this is
clear for $u,v\in A$; one may reduce to this case by multiplying $u$ and $v$
by a common denominator). We obtain more than the existence of a gcd (``the
identity of Bezout''): there exist elements $a,b\in A$ such that the gcd of
$u$ and $v$ may be written in the form
\begin{equation}
d=au+bv.
\label{eq:I.1.2}
\end{equation}
The greatest common divisor of $u$ and $v$ is obviously unique up to
multiplication by units in $A$.

\medskip
\noindent II.\quad Two arbitrary elements $u,v\in K$ have a \emph{least
common multiple} (lcm), i.e.\ there is an element $m\in K$ for which the
relations
\begin{equation}
\text{``}u\mid x\text{ and }v\mid x\text{''}\quad\text{and}\quad
\text{``}m\mid x\text{''}
\label{eq:I.1.3}
\end{equation}
are equivalent. This one can see from the fact that sending $t\mapsto t^{-1}$
in $K$ reverses divisibility. So the proof of the existence of least common
multiples reduces to the proof of the existence of greatest common divisors.
From the relation
\begin{equation}
\operatorname{lcm}(u,v)=\operatorname{gcd}(u^{-1},v^{-1})^{-1}
\quad(u,v\neq0)
\label{eq:I.1.4}
\end{equation}
we easily obtain the following well-known formula
\begin{equation}
\operatorname{gcd}(u,v)\cdot\operatorname{lcm}(u,v)=uv.
\label{eq:I.1.5}
\end{equation}

We may also proceed as in (I) to remark that the existence of
$\operatorname{lcm}(u,v)$ is equivalent to the existence of a greatest lower
bound for $Au$ and $Av$ in the partially ordered set of principal fractional
ideals. The greatest lower bound for $Au$ and $Av$ is $Au\cap Av$.

\medskip
\noindent III.\quad Two elements $a,b$ of $A$ are called \emph{relatively
prime} if $\gcd(a,b)=1$. Let us recall the important \textsc{Lemma of
Euclid}. Let $a,b,c$ be elements in a principal ideal ring $A$. If $a$
divides $bc$ and is relatively prime to $b$, then $a$ divides $c$.

\begin{proof}[A quick proof of Euclid's lemma]
By Bezout \eqref{eq:I.1.2} there exist $a'$ and $b'\in A$ such that
$1=a'a+b'b$; whence $c=a'ac+b'bc$. Since $a$ divides both terms on the
right-hand side, $a$ divides $c$ as well.
\end{proof}

\medskip
\noindent IV.\quad Finally there is unique factorization into products of
primes.

\begin{statement}{Theorem}
Let $A$ be a principal ideal ring and let $K$ be its field of fractions.
There exists a subset $P\subset A$ such that any $x\in K$ may be uniquely
expressed in the form
\begin{equation}
x=u\prod_{p\in P}p^{\nu_{p}(x)},
\label{eq:I.1.6}
\end{equation}
where $u$ is a unit in $A$ and where the exponents $\nu_{p}(x)$ are elements
of $\mathbf{Z}$, all zero except for a finite subset among them.
\end{statement}

For a more systematic exposition of these topics we refer the reader to
\hyperref[bib:1]{[1]}, \emph{Algèbre}, Chapter VI, §1 and Chapter VII, §1. Part of the theory
(more precisely everything which doesn't depend upon Bezout's identity)
extends to more general rings than principal ideal rings. We are referring
to unique factorization domains or factorial rings. See \hyperref[bib:1]{[1]} or \hyperref[bib:2]{[2]},
\emph{Algèbre commutative}, Chapter VII, §3.

%% ----------------------------------------------------------------------
\sechead{2}{An example: the diophantine equations $X^{2}+Y^{2}=Z^{2}$ and
$X^{4}+Y^{4}=Z^{4}$}%
        {AN EXAMPLE: THE DIOPHANTINE EQUATIONS}

One of the most attractive parts of number theory is the study of
\emph{diophantine equations}. One considers polynomial equations
$P(X_{1},\ldots,X_{n})=0$ with coefficients in $\mathbf{Z}$ (respectively,
in $\mathbf{Q}$) and one seeks solutions $(x_{i})$ in $\mathbf{Z}$
(respectively, in $\mathbf{Q}$). One can replace $\mathbf{Z}$
(respectively, $\mathbf{Q}$) by more general rings $A$ (respectively,
fields $K$). We will give an example later (§~6).

We are going to study here two special cases of Fermat's famous equation:
\begin{equation}
X^{n}+Y^{n}=Z^{n}.
\label{eq:I.2.1}
\end{equation}
Fermat claimed to have shown that, for $n\geq3$, this equation has no
non-trivial integer solution $(x,y,z)$. His proof has never been found.
Numerous mathematicians have since Fermat's time worked intensively on this
problem. They have shown that Fermat's claim is true for a great many values
of the exponent $n$. Nevertheless, no general proof (i.e.\ valid for any
$n$) has been found.

Present-day opinion holds that, in his ``proof'', Fermat made a mistake, but
a mistake worthy of a first-class mathematician. For example he might have
conceived the idea (ingenious for his time) of working in the ring of
integers of a field containing $n$th roots of unity; he may have believed
that such a ring is always a principal ideal ring. In fact, we know how to
prove Fermat's claim for any exponent $n$ for which the ring of $n$th roots
of unity is a principal ideal ring. However, this is not the case for all
$n$. For $n$ prime, this ring is a principal ideal ring only for finitely
many values of $n$.\footnote{Cf.\ C.~L. Siegel, \emph{Gesammelte Werke},
Part III, pp.~436--442.}

For $n=2$, equation \eqref{eq:I.2.1} has integer solutions, e.g.\
$(3,4,5)$. One can give a complete description of all the integer solutions
of \eqref{eq:I.2.1}.

\begin{theorem}
\label{thm:I.2.1}
If $x,y,z$ are positive integers such that $x^{2}+y^{2}=z^{2}$, then there
exists an integer $d$ and two relatively prime integers $u$ and $v$ such
that (except, possibly, for a permutation of $x$ and $y$):
\begin{equation}
x=d(u^{2}-v^{2}),\quad y=2duv,\quad\text{and}\quad z=d(u^{2}+v^{2}).
\label{eq:I.2.2}
\end{equation}
\end{theorem}

\begin{proof}
An easy calculation shows that formula \eqref{eq:I.2.2} gives solutions for
$X^{2}+Y^{2}=Z^{2}$. Conversely, let $x,y$, and $z$ be positive integers
such that $x^{2}+y^{2}=z^{2}$. After dividing $x,y,z$ by their greatest
common divisor, we may assume that the three numbers are relatively prime.
It follows that they are pairwise relatively prime as well; for example, if
$x$ and $z$ have a common prime factor $p$, then $p$ divides
$y^{2}=z^{2}-x^{2}$ and, therefore, also $y$. In particular, two of the
numbers $x,y,z$ are odd; the third is necessarily even. The numbers $x$ and
$y$ cannot both be odd, for, if they were, we would have
$x^{2}\equiv1\pmod{4}$, $y^{2}\equiv1\pmod{4}$, and
$z^{2}\equiv2\pmod{4}$, which contradicts the fact that $z^{2}$ is a square.
We have, then, after possibly switching $x$ and $y$:
\begin{equation}
x\text{ odd},\ y\text{ even},\text{ and }z\text{ odd}.
\label{eq:I.2.3}
\end{equation}
Note that
\begin{equation}
y^{2}=z^{2}-x^{2}=(z-x)(z+x).
\label{eq:I.2.4}
\end{equation}
Since the greatest common divisor of $2x$ and $2z$ is $2$, and since
$2x=(z+x)-(z-x)$, $2z=(z+x)+(z-x)$, the greatest common divisor of $z-x$
and $z+x$ can only be $2$. Put $y=2y'$, $z+x=2x'$, $z-x=2z'$, where
$y',x'$, and $z'$ are integers (since $y$, $z+x$, and $z-x$ are even by
\eqref{eq:I.2.3}). We have $y'^{2}=x'z'$. Since $x'$ and $z'$ are relatively
prime, we see that $x'$ and $z'$ are squares $u^{2}$ and $v^{2}$; in fact
any prime factor of $y'^{2}$ appears, with an even exponent, either in the
prime factorization of $x'$ or in that of $z'$, but not in both. We thus
have $z+x=2u^{2}$, $z-x=2v^{2}$, and $y^{2}=2u^{2}\cdot2v^{2}$; whence,
$x=u^{2}-v^{2}$, $y=2uv$, and $z=u^{2}+v^{2}$. Here $u$ and $v$ are
relatively prime, since otherwise $x,y,z$ would have a common prime factor.
Multiplying through by the greatest common divisor of $x,y,z$, call it $d$,
we obtain \eqref{eq:I.2.2}.
\end{proof}

\begin{theorem}
\label{thm:I.2.2}
The equation $X^{4}+Y^{4}=Z^{2}$ has no solution in positive integers
$x,y,z$.
\end{theorem}

\begin{proof}
If there is a solution $(x,y,z)$, where $x,y$, and $z$ are positive
integers, then there is such a solution in which $z$ is minimal. In this
case, $x,y$, and $z$ are pairwise relatively prime; if for example $x$ and
$y$ have a common prime factor $p$, then $p^{4}$ divides $z^{2}$, so
$p^{2}$ divides $z$ and $(x/p,y/p,z/p^{2})$ would be a solution,
contradicting the minimality of $z$. The two other cases are analogous and
even easier. As our equation may be written as $(X^{2})^{2}+(Y^{2})^{2}=Z^{2}$,
we may apply Theorem \ref{thm:I.2.1} to it. After possibly permuting $x$
and $y$ we see that there are positive relatively prime integers $u$ and $v$
such that
\begin{equation}
x^{2}=u^{2}-v^{2},\quad y^{2}=2uv,\quad\text{and}\quad z=u^{2}+v^{2}.
\label{eq:I.2.5}
\end{equation}
Since $4\mid y^{2}$, the relation $y^{2}=2uv$ implies that one of the two
numbers $u$ and $v$ is even; the other is necessarily odd. Thus, $u$ even
and $v$ odd entails $u^{2}\equiv0\pmod{4}$ and $v^{2}\equiv1\pmod{4}$,
whence $x^{2}=u^{2}-v^{2}\equiv-1\pmod{4}$, an impossibility. So $u$ is odd
and $v=2v'$. The relation $y^{2}=4uv'$ and the fact that $u$ and $v'$ are
relatively prime implies that $u$ and $v'$ are squares $a^{2}$ and $b^{2}$.
We apply Theorem \ref{thm:I.2.1} once more, this time to the equation
$x^{2}+v^{2}=u^{2}$ (cf.\ \eqref{eq:I.2.5}). Since $x$ and $u$ are odd, $v$
even, and $x,v$, and $u$ pairwise prime, we obtain two relatively prime
positive integers $c$ and $d$ such that
\begin{equation}
x=c^{2}-d^{2},\quad v=2cd,\quad\text{and}\quad u=c^{2}+d^{2}.
\label{eq:I.2.6}
\end{equation}
Now, from $v=2v'=2b^{2}$, it follows that $cd=b^{2}$, so that $c$ and $d$
are again squares $x'^{2}$ and $y'^{2}$ (they are relatively prime). Since
$u=a^{2}$, \eqref{eq:I.2.6} may be rewritten as
\begin{equation}
a^{2}=x'^{4}+y'^{4},
\label{eq:I.2.7}
\end{equation}
which is of the same form as the original equation. On the other hand, by
\eqref{eq:I.2.5}, $z=u^{2}+v^{2}=a^{4}+4b^{4}>a^{4}$, whence $z>a$, which
contradicts the minimality of $z$. Theorem \ref{thm:I.2.2} is proved.
\end{proof}

A slight variant of our proof shows that, given a solution $(x,y,z)$ in
positive integers of $X^{4}+Y^{4}=Z^{2}$, one may construct an infinite
sequence $(x_{n},y_{n},z_{n})$ of such solutions, where the sequence
$(z_{n})$ is strictly decreasing. This is an absurdity. This method of proof
is called the \emph{method of infinite descent} and is due to Fermat.

\begin{statement}{Corollary}
The equation $X^{4}+Y^{4}=Z^{4}$ has no positive integer solutions.
\end{statement}

\begin{proof}
This equation may be written in the form $X^{4}+Y^{4}=(Z^{2})^{2}$, to which
Theorem \ref{thm:I.2.2} applies.
\end{proof}

%% ----------------------------------------------------------------------
\sechead{3}{Some lemmas concerning ideals; Euler's
$\varphi$-function}%
        {SOME LEMMAS CONCERNING IDEALS; EULER'S $\varphi$-FUNCTION}

Let $n\geq1$ be a natural number. We write $\varphi(n)$ for the number of
integers $q$, $0\leq q\leq n$, such that $q$ and $n$ are relatively prime
(since $0$ and $n$ are divisible by $n$, it is equivalent to take
$1\leq q\leq n-1$ for any $n>1$; set $\varphi(1)=1$). The function $\varphi$
so defined is called \emph{Euler's $\varphi$-function}. If $p$ is a prime
number, then clearly:
\begin{equation}
\varphi(p)=p-1.
\label{eq:I.3.1}
\end{equation}
For $n=p^{s}$, a power of a prime, the integers relatively prime to $n$ are
those integers which are not multiples of $p$. There are $p^{s}-1$ multiples
of $p$ between $1$ and $p^{s}$. Therefore,
\begin{equation}
\varphi(p^{s})=p^{s}-p^{s-1}=p^{s-1}(p-1).
\label{eq:I.3.2}
\end{equation}

Now we intend to calculate $\varphi(n)$ by making use of the fact that $n$
may be expressed as a product of powers of primes. For this purpose we need
some properties of $\varphi(n)$ and we need some lemmas concerning ideals.
These lemmas will be useful later.

\begin{proposition}
\label{prop:I.3.1}
Let $n\geq1$ be a natural number. The value $\varphi(n)$ of Euler's
$\varphi$-function equals the number of elements of $\mathbf{Z}/n\mathbf{Z}$
which generate this group. It also equals the number of units in the ring
$\mathbf{Z}/n\mathbf{Z}$.
\end{proposition}

\begin{proof}
Let us recall that each congruence class mod $n\mathbf{Z}$ contains a unique
integer $q$ such that $0\leq q\leq n-1$. For such an integer $q$ we write
$\bar{q}$ for its residue class mod $n\mathbf{Z}$. It suffices to prove the
following implications: $q$ relatively prime to $n\Rightarrow\bar{q}$ a unit
in the ring $\mathbf{Z}/n\mathbf{Z}\Rightarrow\bar{q}$ generates the
additive group $\mathbf{Z}/n\mathbf{Z}\Rightarrow q$ relatively prime to
$n$. If $q$ is relatively prime to $n$, Bezout's identity
(\hyperref[sec:1.1]{§1}, \eqref{eq:I.1.2}) implies that there are integers $x$ and $y$ such that
$qx+ny=1$; whence $\bar{q}\,\bar{x}=\overline{1}$, so $\bar{q}$ is a unit in
$\mathbf{Z}/n\mathbf{Z}$. Let $\bar{q}$ be a unit in
$\mathbf{Z}/n\mathbf{Z}$. Writing $x$ for an integer such that
$\bar{q}\cdot\bar{x}=1$, we see that
$\bar{a}=\bar{a}\cdot\bar{x}\,\bar{q}$ (in the ring
$\mathbf{Z}/n\mathbf{Z}$), where $\bar{a}$ is an arbitrary element of
$\mathbf{Z}/n\mathbf{Z}$ ($0\leq a<n$). It follows that
$\bar{a}=(ax)\cdot\bar{q}$ (in the additive group
$\mathbf{Z}/n\mathbf{Z}$), so $\bar{q}$ generates the group
$\mathbf{Z}/n\mathbf{Z}$. Finally, if $\bar{q}$ generates
$\mathbf{Z}/n\mathbf{Z}$, there is an $x$ such that $x\bar{q}=\overline{1}$,
thus such that $xq\equiv1\pmod{n}$; thus there exists an integer $y$ for
which $xq-1=yn$, so $1=xq-yn$. This is an instance of Bezout's identity,
which shows that $q$ is relatively prime to $n$.
\end{proof}

\begin{lemma}
\label{lem:I.3.1}
Let $A$ be a ring, $\mathfrak{a}$ and $\mathfrak{b}$ ideals of $A$ such that
$\mathfrak{a}+\mathfrak{b}=A$. Then
$\mathfrak{a}\cap\mathfrak{b}=\mathfrak{a}\mathfrak{b}$ and the canonical
homomorphism $\varphi:A\to A/\mathfrak{a}\times A/\mathfrak{b}$ induces an
isomorphism $\theta:A/\mathfrak{a}\mathfrak{b}\to
A/\mathfrak{a}\times A/\mathfrak{b}$.
\end{lemma}

Recall that the homomorphism $\varphi$ sends any $x\in A$ into the pair
consisting of the class of $x\bmod\mathfrak{a}$ and the class of
$x\bmod\mathfrak{b}$.

\begin{proof}
We know that, in general, $\mathfrak{a}\mathfrak{b}\subset\mathfrak{a}$ and
$\mathfrak{a}\mathfrak{b}\subset\mathfrak{b}$, so
$\mathfrak{a}\mathfrak{b}\subset\mathfrak{a}\cap\mathfrak{b}$. Let
$x\in\mathfrak{a}\cap\mathfrak{b}$. Since $\mathfrak{a}+\mathfrak{b}=A$,
there are elements $a\in\mathfrak{a}$ and $b\in\mathfrak{b}$ such that
$a+b=1$. Thus $x=ax+xb$ is a sum of two elements of
$\mathfrak{a}\mathfrak{b}$, whence $x\in\mathfrak{a}\mathfrak{b}$ and
$\mathfrak{a}\cap\mathfrak{b}\subset\mathfrak{a}\mathfrak{b}$. Therefore
$\mathfrak{a}\mathfrak{b}=\mathfrak{a}\cap\mathfrak{b}$.

Clearly $\mathfrak{a}\cap\mathfrak{b}$ is the kernel of $\varphi$. Since
$\mathfrak{a}\cap\mathfrak{b}=\mathfrak{a}\mathfrak{b}$, $\varphi$ is
constant on each residue class mod $\mathfrak{a}\mathfrak{b}$. Thus, we
obtain a mapping
$\theta:A/\mathfrak{a}\mathfrak{b}\to A/\mathfrak{a}\times A/\mathfrak{b}$,
which is obviously a homomorphism. Since
$\varphi^{-1}(0)=\mathfrak{a}\mathfrak{b}$, $\theta^{-1}(0)=(0)$, so $\theta$
is injective. It remains to show that $\theta$ is surjective.

We have ``passed to the quotient ring'' in the course of the above
argument. Henceforth, we shall be much more brief in explaining analogous
arguments.

In order to show that $\theta$ (or, equivalently, $\varphi$) is surjective,
we have to find, for any pair $y\in A$ and $z\in A$, an element $x\in A$
such that $x+\mathfrak{a}=y+\mathfrak{a}$ and
$x+\mathfrak{b}=z+\mathfrak{b}$. Take $a\in\mathfrak{a}$ and
$b\in\mathfrak{b}$ such that $a+b=1$. Set $x=az+by$. Modulo
$\mathfrak{a}$, $x\equiv by\equiv(1-a)y\equiv y-ay\equiv y$; similarly,
$x\equiv z\bmod\mathfrak{b}$.
\end{proof}

\begin{lemma}
\label{lem:I.3.2}
Let $A$ be a ring and $(\mathfrak{a}_{i})_{1\leq i\leq r}$, a finite set of
ideals of $A$ such that
$\mathfrak{a}_{i}+\mathfrak{a}_{j}=A$ for $i\neq j$. Then there is a
canonical isomorphism of $A/\mathfrak{a}_{1}\ldots\mathfrak{a}_{r}$ onto
$\prod_{i=1}^{r}A/\mathfrak{a}_{i}$.
\end{lemma}

\begin{proof}
Lemma \ref{lem:I.3.1} is the case $r=2$ of Lemma \ref{lem:I.3.2}. We
proceed by induction on $r$. Put
$\mathfrak{b}=\mathfrak{a}_{2}\ldots\mathfrak{a}_{r}$. Let us show that
$\mathfrak{a}_{1}+\mathfrak{b}=A$. For $i\geq2$ we have
$\mathfrak{a}_{1}+\mathfrak{a}_{i}=A$, so there are elements
$c_{i}\in\mathfrak{a}_{1}$ and $a_{i}\in\mathfrak{a}_{i}$ such that
\[
c_{i}+a_{i}=1,\qquad
1=\prod_{i=1}^{r}(c_{i}+a_{i})=c+a_{2}\ldots a_{r},
\]
where $c$ is a sum of terms each of which contains at least one $c_{i}$ as a
factor. Therefore, $c\in\mathfrak{a}_{1}$. As
$\mathfrak{a}_{2}\ldots\mathfrak{a}_{r}\in\mathfrak{b}$, it follows that
$\mathfrak{a}_{1}+\mathfrak{b}=A$.

By Lemma \ref{lem:I.3.1}, $A/\mathfrak{a}_{1}\mathfrak{b}$ is isomorphic to
$A/\mathfrak{a}_{1}\times A/\mathfrak{b}$. According to the induction
hypothesis $A/\mathfrak{b}=A/\mathfrak{a}_{2}\ldots\mathfrak{a}_{r}$ which
is isomorphic to $A/\mathfrak{a}_{2}\times\cdots\times A/\mathfrak{a}_{r}$.
The lemma follows by composing these isomorphisms.
\end{proof}

Let us apply these lemmas to $\mathbf{Z}$.

\begin{proposition}
\label{prop:I.3.2}
Let $n$ and $n'$ be relatively prime integers. Then the ring
$\mathbf{Z}/nn'\mathbf{Z}$ is isomorphic to the product ring
$\mathbf{Z}/n\mathbf{Z}\times\mathbf{Z}/n'\mathbf{Z}$.
\end{proposition}

\begin{proof}
This is a special case of Lemma \ref{lem:I.3.1}, the hypothesis
$n\mathbf{Z}+n'\mathbf{Z}=\mathbf{Z}$ being Bezout's identity.
\end{proof}

\begin{corollary}
\label{cor:I.3.1}
If $n$ and $n'$ are relatively prime positive integers, then
$\varphi(nn')=\varphi(n)\varphi(n')$.
\end{corollary}

\begin{proof}
$\varphi(nn')$ is the number of units in $\mathbf{Z}/nn'\mathbf{Z}$
(Proposition \ref{prop:I.3.1}), which is isomorphic to
$\mathbf{Z}/n\mathbf{Z}\times\mathbf{Z}/n'\mathbf{Z}$. Now an element
$(\alpha,\beta)$ of a ring product is invertible if and only if each of its
components $\alpha,\beta$ is invertible. Thus our assertion follows from
Proposition \ref{prop:I.3.1}.
\end{proof}

\begin{corollary}
\label{cor:I.3.2}
Let $n$ be a positive integer and let
$n=p_{1}^{\alpha_{1}}\cdots p_{r}^{\alpha_{r}}$ be its prime factorization.
Then $\varphi(n)=n(1-1/p_{1})\cdots(1-1/p_{r})$.
\end{corollary}

\begin{proof}
By Corollary \ref{cor:I.3.1},
$\varphi(n)=\varphi(p_{1}^{\alpha_{1}})\cdots
\varphi(p_{r}^{\alpha_{r}})$. By \eqref{eq:I.3.2},
$\varphi(p_{i}^{\alpha_{i}})=p_{i}^{\alpha_{i}-1}(p_{i}-1)
=p_{i}^{\alpha_{i}}(1-1/p_{i})$. Multiplication gives our formula.
\end{proof}

%% ----------------------------------------------------------------------
\sechead{4}{Some preliminaries concerning modules}%
        {SOME PRELIMINARIES CONCERNING MODULES}

Before studying modules over a principal ideal ring, we make some remarks
concerning modules over arbitrary commutative rings.

Let $A$ be a commutative ring and let $I$ be a set. Let $A^{(I)}$ denote the
set of sequences $(a_{i})_{i\in I}$, indexed by $I$, of elements of $A$ such
that $a_{i}=0$ except for a finite number of indices $i\in I$. Thus
$A^{(I)}$ is a subset of the cartesian product set $A^{I}$, and also a
submodule of $A^{I}$ if one provides $A^{I}$ with the $A$-module structure
defined by componentwise addition and scalar multiplication.

If $I$ is finite, then $A^{(I)}=A^{I}$.

For $j\in I$, the sequence $(\delta_{ij})_{i\in I}$ such that
$\delta_{jj}=1$ and $\delta_{ij}=0$ for $i\neq j$ is an element $e_{j}$ of
$A^{(I)}$. Every element $(a_{j})_{j\in I}$ of $A^{(I)}$ has a unique
expression as a (finite) linear combination of the $e_{j}$. More precisely,
\begin{equation}
(a_{j})_{j\in I}=\sum_{j\in I}a_{j}e_{j}
\label{eq:I.4.1}
\end{equation}
(note that, in the summation on the right, all the terms are zero except for
a finite number, so that the summation makes sense).

We call $(e_{j})_{j\in I}$ the \emph{canonical base} for $A^{(I)}$.

Let $A$ be a ring, $M$ an $A$-module, and $(x_{i})_{i\in I}$ a family of
elements of $M$. To every element $(a_{i})_{i\in I}$ of $A^{(I)}$ let us
associate the element $\sum_{i}a_{i}x_{i}$ of $M$ (as before, the summation
makes sense). Thus we obtain a mapping $\varphi:A^{(I)}\to M$, which is
obviously linear. If $(e_{i})_{i\in I}$ is the canonical base for $A^{(I)}$,
then $\varphi(e_{i})=x_{i}$ for any $i\in I$. The equivalence of the
following statements is immediate:
\begin{equation}
(x_{i})_{i\in I}\text{ is a linearly independent set if and only if }
\varphi\text{ is injective;}
\label{eq:I.4.2}
\end{equation}
\begin{equation}
(x_{i})_{i\in I}\text{ generates }M\text{ if and only if }\varphi
\text{ is surjective.}
\label{eq:I.4.3}
\end{equation}
If $\varphi$ is bijective $(x_{i})_{i\in I}$ is called a \emph{base} for
$M$. This means that every element of $M$ has a unique expression as a
linear combination of the elements $(x_{i})_{i\in I}$. A module $M$ which
has a base is called a \emph{free module}.

In contrast to the case of vector spaces over a field, a module over a ring
does not necessarily have a base. For example, the $\mathbf{Z}$-module
$\mathbf{Z}/n\mathbf{Z}$ for $n\neq0$ or $1$. In the ensuing discussion we
shall show that certain modules are free. This type of result is seldom
trivial.

A module is said to be \emph{of finite type} if it contains a finite
generating set. The following theorem is basic for the study of the
properties of Noetherian rings and modules. We develop this topic further in
\hyperref[chap:III]{Chapter III}.

\begin{theorem}
\label{thm:I.4.1}
Let $A$ be a ring and $M$ an $A$-module. The following conditions are
equivalent.
\begin{alist}
\item Every non-empty family of submodules of $M$ contains a maximal element
(under the relation of inclusion).
\item Every increasing sequence $(M_{n})_{n\geq0}$ (again for the relation
of inclusion) of submodules of $M$ is stationary (i.e.\ there exists $n_{0}$
such that $M_{n}=M_{n_{0}}$ for all $n\geq n_{0}$).
\item Every submodule of $M$ is of finite type.
\end{alist}
\end{theorem}

\begin{proof}
Let us show that (a) implies (c). Let $E$ be a submodule of $M$ and let
$\Phi$ be the collection consisting of all submodules of finite type of $E$.
Then $\Phi$ is not empty, since $(0)\in\Phi$. By (a) $\Phi$ contains a
maximal element $F$. For $x\in E$, $F+Ax$ is a submodule of finite type of
$E$ (it is generated by the union of $\{x\}$ and any finite set of
generators for $F$). Thus $F+Ax=F$, since $F+Ax\supset F$ and $F$ is
maximal. Therefore, $x\in F$, $E\subset F$, $E=F$, and $E$ is of finite
type.

We prove now that (c) implies (b). Let $(M_{n})_{n\geq0}$ be an increasing
sequence of submodules of $M$. Then $E=\bigcup_{n\geq0}M_{n}$ is a
submodule of $M$. By (c) it contains a finite set of generators
$(x_{1},\ldots,x_{q})$. For every $i$ there is an index $n(i)$ such that
$x_{i}\in M_{n(i)}$. Let $n_{0}$ be the largest of the $n(i)$'s. Then
$x_{i}\in M_{n_{0}}$ for all $i$, so $E\subset M_{n_{0}}$ and
$E=M_{n_{0}}$. For $n\geq n_{0}$ the inclusion relations
$M_{n_{0}}\subset M_{n}\subset E$ and the equality $M_{n_{0}}=E$ imply that
$M_{n_{0}}=M_{n}$. Thus the sequence $(M_{n})$ is stationary beyond $n_{0}$.

It remains to prove that (b) implies (a). The equivalence of (a) and (b) is
a special case of the following lemma concerning partially ordered sets.
\end{proof}

\begin{lemma}
\label{lem:I.4.1}
Let $T$ be a partially ordered set. The following statements are equivalent:
\begin{alist}
\item Every non-empty subset of $T$ contains a maximal element.
\item Every increasing sequence $(t_{n})_{n\geq0}$ of elements of $T$ is
stationary.
\end{alist}
\end{lemma}

\begin{proof}
(a) $\Rightarrow$ (b): Let $t_{q}$ be a maximal element of the increasing
sequence $(t_{n})$. Then, for $n\geq q$, $t_{n}\geq t_{q}$ (the sequence
increases), so $t_{n}=t_{q}$ (maximality).

(b) $\Rightarrow$ (a): Suppose there exists a subset $S$ of $T$ which does
not contain a maximal element. Then, for any $x\in S$, the set of elements
of $S$ which are larger than $x$ is nonempty. By the axiom of choice there
exists a mapping $f:S\to S$ such that $f(x)>x$ for all $x\in S$. Since $S$
is not empty, one may choose $t_{0}\in S$ and define by induction the
sequence $(t_{n})_{n\geq0}$ by setting $t_{n+1}=f(t_{n})$. This sequence is
strictly increasing; it is therefore not stationary. This contradicts (b),
so (b) $\Rightarrow$ (a) is established.
\end{proof}

\begin{statement}{Corollary of theorem \ref{thm:I.4.1}}
In a principal ideal ring $A$, every non-empty family of ideals contains a
maximal element.
\end{statement}

\begin{proof}
If one considers $A$ as a module over itself, its submodules are its ideals.
As all ideals are principal, they are $A$-modules generated by a single
element, thus of finite type. The corollary follows from the implication
(c) $\Rightarrow$ (a) of Theorem \ref{thm:I.4.1}.
\end{proof}

%% ----------------------------------------------------------------------
\sechead{5}{Modules over principal ideal rings}%
        {MODULES OVER PRINCIPAL IDEAL RINGS}

Let $A$ be an integral domain and let $K$ be its field of fractions. A free
$A$-module (isomorphic to an $A^{(I)}$ for some $I$) may be injected into a
vector space over $K$ ($K^{(I)}$ in the case of $A^{(I)}$). It follows that
the same thing is true for any submodule $M$ of a free $A$-module. The
dimension of the subspace generated by $M$ is called the \emph{rank} of $M$.
If $M$ is itself free and admits a base having $n$ elements, then the rank
of $M$ is $n$.

\begin{theorem}
\label{thm:I.5.1}
Let $A$ be a principal ideal ring, $M$ a free $A$-module of rank $n$, and
$M'$ a submodule of $M$. Then:
\begin{alist}
\item $M'$ is free of rank $q$, $0\leq q\leq n$.
\item If $M'\neq(0)$, there exists a base $(e_{1},\ldots,e_{n})$ of $M$ and
non-zero elements $a_{1},\ldots,a_{q}\in A$ such that
$(a_{1}e_{1},\ldots,a_{q}e_{q})$ is a base of $M'$ and such that $a_{i}$
divides $a_{i+1}$, $1\leq i\leq q-1$.
\end{alist}
\end{theorem}

\begin{proof}
The theorem is trivial for $M'=(0)$, so we may assume $M'\neq(0)$. Let
$L(M,A)$ be the set of linear forms on $M$. For $u\in L(M,A)$, $u(M')$ is a
submodule of $A$, an ideal of $A$. We may write $u(M')=Aa_{u}$ with
$a_{u}\in A$, since the ideal is principal. Let $u\in L(M,A)$ be such that
$Aa_{u}$ is maximal among the $Aa_{v}$ ($v\in L(M,A)$) (\hyperref[sec:1.4]{§4}, corollary of
Theorem \ref{thm:I.4.1}). Let us take a base $(x_{1},\ldots,x_{n})$ of $M$,
which identifies $M$ with $A^{n}$. Let $pr_{i}:M\to A$ be the projection on
the $i$th coordinate, i.e.\ $pr_{i}(x_{j})=\delta_{ij}$. Since
$M'\neq(0)$, for at least one $i$, $1\leq i\leq n$, $pr_{i}(M')$ is not
$(0)$. Thus $a_{u}\neq(0)$. By our construction there exists $e'\in M'$ such
that $u(e')=a_{u}$. Let us show that for every $v\in L(M,A)$,
$a_{u}\mid v(e')$. Indeed, if $d=\gcd(a_{u},v(e'))$, then
$d=ba_{u}+cv(e')$ with $b,c\in A$, whence $d=(bu+cv)(e')$. Since $bu+cv$ is
a linear form on $M$, $Aa_{u}\subset Ad\subset u(M')$. The maximality of
$Aa_{u}$ implies that $Ad=Aa_{u}$, so $a_{u}$ must divide $v(e')$.

In particular, $a_{u}\mid pr_{i}(e')$, so let $pr_{i}(e')=a_{u}b_{i}$ with
$b_{i}\in A$. Put $e=\sum_{i=1}^{n}b_{i}x_{i}$. Then $e'=a_{u}e$. Since
$u(e')=a_{u}=a_{u}\cdot u(e)$, it follows that $u(e)=1$ (note that
$a_{u}\neq0$). Let us show that
\begin{equation}
M=\operatorname{Ker}(u)+Ae
\label{eq:I.5.1}
\end{equation}
and
\begin{equation}
M'=(M'\cap\operatorname{Ker}(u))+Ae'
\label{eq:I.5.2}
\end{equation}
(where $e'=a_{u}e$), the sums being direct.
\end{proof}

\noindent 1.\quad For every $x\in M$,
$x=u(x)e+(x-u(x)e)$. We see that
$u(x-u(x)e)=u(x)-u(x)u(e)=0$, since $u(e)=1$, so
$x-u(x)e\in\operatorname{Ker}(u)$. This shows that
$Ae+\operatorname{Ker}(u)=M$; obviously
$Ae\cap\operatorname{Ker}(u)=(0)$.

\noindent 2.\quad For $y\in M'$, $u(y)=ba_{u}$ with $b\in A$, so
$y=ba_{u}e+(y-u(y)e)=be'+(y-u(y)e)$.

Again it is clear that $y-u(y)e\in\operatorname{Ker}(u)$ and also that
$y-u(y)e=y-be'\in M'$, i.e.\
$y-u(y)e\in M'\cap\operatorname{Ker}(u)$ and $be'\in Ae'\subset Ae$. This
entails \eqref{eq:I.5.2}.

Now we prove (a) by induction on the rank $q$ of $M'$. If $q=0$, $M'=(0)$
and there is nothing to prove. If $q>0$,
$M'\cap\operatorname{Ker}(u)$ is of rank $q-1$ according to
\eqref{eq:I.5.2}, and is therefore free according to the induction
hypothesis. As, in \eqref{eq:I.5.2}, the sum is direct, we obtain a base for
$M'$ by adding $e'$ to a base for $M'\cap\operatorname{Ker}(u)$. Thus $M'$
is free and (a) is true.

To prove (b) we argue by induction on the rank $n$ of $M$. Again the case
$n=0$ is trivial. By (a), $\operatorname{Ker}(u)$ is free and of rank
$n-1$, since, in \eqref{eq:I.5.1}, the sum is direct. We apply the induction
hypothesis to the free module $\operatorname{Ker}(u)$ and to its submodule
$M'\cap\operatorname{Ker}(u)$: if
$M'\cap\operatorname{Ker}(u)\neq(0)$, there exists $q\leq n$, a base
$(e_{2},\ldots,e_{n})$ of $\operatorname{Ker}(u)$, and there are non-zero
elements $a_{2},\ldots,a_{q}$ of $A$ such that
$(a_{2}e_{2},\ldots,a_{q}e_{q})$ is a base for
$M'\cap\operatorname{Ker}(u)$ and such that $a_{i}$ divides $a_{i+1}$,
$2\leq i\leq q-1$. Keeping the same notations as above, we set $a_{1}=a_{u}$
and $e_{1}=e$. Then $(e_{1},e_{2},\ldots,e_{n})$ is a base for $M$ (according
to \eqref{eq:I.5.1}), and $(a_{1}e_{1},\ldots,a_{q}e_{q})$ is a base for
$M'$ (according to \eqref{eq:I.5.2} and the fact that $e'=a_{1}e_{1}$). It
remains to prove that $a_{1}\mid a_{2}$. Let $v$ be the linear form on $M$
defined by the relation $v(e_{1})=v(e_{2})=1$ and $v(e_{i})=0$ for
$i\geq3$. Then $a_{1}=a_{u}=v(a_{u}e_{1})=v(e')\in v(M')$, so
$Aa_{u}\subset v(M')$. By the maximality of $Aa_{u}$ we may conclude that
$v(M')=Aa_{u}=Aa_{1}$. Since $a_{2}=v(a_{2}e_{2})\in v(M')$, we see that
$a_{2}\in Aa_{1}$, i.e.\ $a_{1}\mid a_{2}$.

The ideals $Aa_{i}$ of Theorem \ref{thm:I.5.1} are called the
\emph{invariant factors} of $M'$ in $M$. One can show that they are uniquely
determined by $M$ and $M'$ (\hyperref[bib:1]{[1]}, Chapter VII, §3).

\begin{corollary}
\label{cor:I.5.1}
Let $A$ be a principal ideal ring. Let $E$ be an $A$-module of finite type.
Then $E$ is isomorphic to a product
$(A/\mathfrak{a}_{1})\times(A/\mathfrak{a}_{2})\times\cdots\times
(A/\mathfrak{a}_{n})$, where the $\mathfrak{a}$'s are ideals of $A$ such that
$\mathfrak{a}_{1}\supset\mathfrak{a}_{2}\supset\cdots\supset
\mathfrak{a}_{n}$.
\end{corollary}

\begin{proof}
Let $(x_{1},\ldots,x_{n})$ be a finite set of generators of $E$. According
to the beginning of \hyperref[sec:1.4]{§4}, there is a surjective homomorphism
$\varphi:A^{n}\to E$, such that $E$ is isomorphic to
$A^{n}/\operatorname{Ker}(\varphi)$. By Theorem \ref{thm:I.5.1}, there is a
base $(e_{1},\ldots,e_{n})$ of $A^{n}$, an integer $q\leq n$, and nonzero
elements $a_{1},\ldots,a_{q}$ of $A$ such that
$(a_{1}e_{1},\ldots,a_{q}e_{q})$ is a base of $\operatorname{Ker}(\varphi)$
and such that $a_{i}$ divides $a_{i+1}$ for all $i$, $1\leq i\leq q-1$. Put
$a_{p}=0$ for $q+1\leq p\leq n$. Then $A^{n}/\ker(\varphi)$ is isomorphic to
the product of the $Ae_{i}/Aa_{i}e_{i}$ ($1\leq i\leq n$), and
$Ae_{i}/Aa_{i}e_{i}$ is isomorphic to
$A/Aa_{i}$. Putting $\mathfrak{a}_{i}=Aa_{i}$, we obtain the corollary.
\end{proof}

We shall say that a module $E$ over an integral domain $A$ is torsion free
if the relation $ax=0$ ($a\in A$, $x\in E$) implies $a=0$ or $x=0$.

\begin{corollary}
\label{cor:I.5.2}
Over a principal ideal ring $A$, every module $E$ of finite type which is
torsion free is free.
\end{corollary}

\begin{proof}
We make use of Corollary \ref{cor:I.5.1}:
$E\simeq(A/\mathfrak{a}_{1})\times\cdots\times(A/\mathfrak{a}_{n})$.
Suppressing the factors which are zero, we may suppose that
$\mathfrak{a}_{i}\neq A$ for all $i$. If $\mathfrak{a}_{1}\neq(0)$, if $a$ is
a non-zero element of $\mathfrak{a}_{1}$, if $x_{1}$ is a non-zero element
of $A/\mathfrak{a}_{1}$, and if $x=(x_{1},0,\ldots,0)$, then $ax=0$ ---
contradicting the fact that $E$ is torsion free. Thus
$\mathfrak{a}_{1}=(0)$, $\mathfrak{a}_{i}=(0)$ for all $i$ (since
$\mathfrak{a}_{i}\subset\mathfrak{a}_{1}$), and $E$ is isomorphic to $A^{n}$.
\end{proof}

The hypothesis that $E$ is of finite type is essential: for example
$\mathbf{Q}$ is a torsion free $\mathbf{Z}$-module which is not free.

\begin{corollary}
\label{cor:I.5.3}
Over a principal ideal ring $A$, every module $E$ of finite type is
isomorphic to a finite product of modules $M_{i}$, where each $M_{i}$ is
equal to $A$ or to a quotient $A/Ap^{s}$ with $p$ prime.
\end{corollary}

\begin{proof}
We make use of Corollary \ref{cor:I.5.1} and we decompose each factor
$A/Aa$, where $a\neq0$, by means of \hyperref[sec:1.3]{§3}, Lemma \ref{lem:I.3.2}: if
$a=up_{1}^{s_{1}}\ldots p_{r}^{s_{r}}$ is the prime factorization of $a$,
$A/Aa$ is isomorphic to the product of the $A/Ap_{i}^{s_{i}}$'s.
\end{proof}

\begin{corollary}
\label{cor:I.5.4}
Let $G$ be a finite commutative group. There exists $x\in G$ whose order is
the least common multiple of the orders of the elements of $G$.
\end{corollary}

\begin{proof}
A commutative group is a $\mathbf{Z}$-module (the operation being addition).
According to Corollary \ref{cor:I.5.1},
\[
G\simeq\mathbf{Z}/a_{1}\mathbf{Z}\times\cdots\times
\mathbf{Z}/a_{n}\mathbf{Z},
\]
where $a_{1}\mid a_{2}\mid\ldots\mid a_{n}$. We have $a_{i}\neq0$ for all
$i$; otherwise $G$ would be infinite. We write $y$ for the residue class of
$1$ in $\mathbf{Z}/a_{n}\mathbf{Z}$ and we put $x=(0,\ldots,0,y)$. The order
of $x$ is obviously $a_{n}$. For $z=(z_{1},\ldots,z_{n})\in G$, we have
$a_{n}z=0$, since $a_{i}$ divides $a_{n}$ for all $i$. Therefore $a_{n}$ is
a multiple of the order of $z$ and $x$ is the element sought.
\end{proof}

%% ----------------------------------------------------------------------
\sechead{6}{Roots of unity in a field}{ROOTS OF UNITY IN A FIELD}

\begin{theorem}
\label{thm:I.6.1}
Let $K$ be a field. Every finite subgroup $G$ of the multiplicative group
$K^{*}$ consists of roots of unity and is cyclic.
\end{theorem}

\begin{proof}
According to \hyperref[sec:1.5]{§5}, Corollary \ref{cor:I.5.4} of Theorem \ref{thm:I.5.1},
there exists $z\in G$ whose order $n$ is such that $y^{n}=1$ for every
$y\in G$. Since a polynomial of degree $n$ over a field (for example
$X^{n}-1$) has at most $n$ roots in the field, the number of elements of
$G$ is at most $n$. Now, inasmuch as $z$ has order $n$, $G$ contains the
$n$ elements $z,z^{2},\ldots,z^{n}=1$, which are all distinct. Thus $G$ is
comprised of these elements and is cyclic.
\end{proof}

If a field $K$ contains $n$ $n$th roots of unity, they form a cyclic group of
order $n$ (isomorphic to $\mathbf{Z}/n\mathbf{Z}$). A generator of this
group is called a \emph{primitive $n$th root of unity}, every $n$th root of
unity is thus a power of such a primitive root. According to \hyperref[sec:1.3]{§3},
Proposition \ref{prop:I.3.1}, the number of these primitive roots is
$\varphi(n)$.

%% ----------------------------------------------------------------------
\sechead{7}{Finite fields}{FINITE FIELDS}

Let $K$ be a field. There is a unique ring homomorphism $\varphi:
\mathbf{Z}\to K$ (defined by $\varphi(n)=1+1+\cdots+1$, $n$ times, for
$n\geq0$ and by $\varphi(-n)=-\varphi(n)$). If $\varphi$ is injective, it
identifies $\mathbf{Z}$ with a subring of $K$; then $K$ also contains the
field of fractions $\mathbf{Q}$ of $\mathbf{Z}$. In this case we say that
$K$ is of \emph{characteristic} $0$. If $\varphi$ is not injective, its
kernel is an ideal $p\mathbf{Z}$ where $p>0$; then $\mathbf{Z}/p\mathbf{Z}$
is identified with a subring of $K$; thus $\mathbf{Z}/p\mathbf{Z}$ is an
integral domain (in fact, a field) from which it follows that $p$ is a prime
number. We say, in this case, that $K$ is of \emph{characteristic} $p$.
From here on, we write $\mathbf{F}_{p}$ for $\mathbf{Z}/p\mathbf{Z}$.

The subfield, $\mathbf{Q}$ or $\mathbf{F}_{p}$, of $K$ is the smallest
subfield of $K$; it is called the \emph{prime subfield} of $K$. For every
prime number $p$ there exist fields of characteristic $p$, e.g.\
$\mathbf{F}_{p}$.

\begin{proposition}
\label{prop:I.7.1}
If $K$ is a field of characteristic $p\neq0$, then $px=0$ for every $x\in K$
and $(x+y)^{p}=x^{p}+y^{p}$ for every $x,y\in K$.
\end{proposition}

\begin{proof}
For $x\in K$, we have $p\cdot x=(p\cdot1)\cdot x=0\cdot x=0$. On the other
hand, the binomial formula gives
\[
(x+y)^{p}=x^{p}+y^{p}+\sum_{j=1}^{p-1}\binom{p}{j}x^{j}y^{p-j}.
\]
The binomial coefficient $\binom{p}{j}$ is an integer; its value is
$p!/[j!(p-j)!]$. Inasmuch as the prime $p$ appears in the numerator but not
in the denominator, $\binom{p}{j}$ is a multiple of $p$ for
$1\leq j\leq p-1$. The intermediate terms in the expansion of $(x+y)^{p}$
thus vanish in a field of characteristic $p$.
\end{proof}

By induction one sees that $(x+y)^{p^{n}}=x^{p^{n}}+y^{p^{n}}$ for every
$n\geq0$.

\begin{theorem}
\label{thm:I.7.1}
Let $K$ be a finite field. Set $q=\operatorname{card}(K)$. Then:
\begin{alist}
\item The characteristic of $K$ is a prime $p$, $K$ is a finite dimensional
vector space of dimension $s$ over $\mathbf{F}_{p}$, and $q=p^{s}$.
\item The multiplicative group $K^{*}$ is cyclic of order $q-1$.
\item $x^{q-1}=1$ for every $x\in K^{*}$; $x^{q}=x$ for every $x\in K$.
\end{alist}
\end{theorem}

\begin{proof}
Since $\mathbf{Z}$ is infinite, $K$ cannot be of characteristic zero. Thus
$K$ contains $\mathbf{F}_{p}$, $p$ prime; in fact $K$ is a vector space over
$\mathbf{F}_{p}$, whose dimension $s$ must be finite---otherwise $K$ would
be an infinite field. As a vector space, $K$ is isomorphic to
$(\mathbf{F}_{p})^{s}$, so $K$ contains $p^{s}$ elements. We see that (b)
follows from \hyperref[sec:1.6]{§6}, Theorem \ref{thm:I.6.1}, and that (c) is an immediate
consequence of (b).
\end{proof}

\begin{statement}{Example}
Let us interpret (b) for the case of $\mathbf{F}_{p}$, $p$ prime. There
exists an integer $x\in\mathbf{Z}$ such that $0<x\leq p-1$ and such that
every integer $y$ which is not a multiple of $p$ is congruent modulo $p$ to
a power of $x$. Such an $x$ is called a \emph{primitive root modulo} $p$.
The problem of finding primitive roots modulo $p$ is by no means trivial.
For instance there are $\varphi(6)=2$ roots primitive modulo $7$; they are
$3$ and $5$ (one sees that $1^{2}\equiv6^{2}\equiv1\pmod{7}$ and
$2^{3}\equiv4^{3}\equiv1\pmod{7}$, $3$ and $5$ are the only other
possibilities).
\end{statement}

\begin{statement}{Remark}
(a) and (c) imply that a finite field $K$ with $q$ elements is the set of
roots of the polynomial $X^{q}-X$ (which has exactly $q$ roots). One can
show that two finite fields with $q$ elements are isomorphic. We write
$\mathbf{F}_{q}$ unambiguously for a field with $q$ elements.
\end{statement}

As an exercise we are going to digress in order to prove the following
elegant theorem which concerns diophantine equations over a finite field.

\begin{theorem}[Chevalley]
\label{thm:I.7.2}
Let $K$ be a finite field and let $F(X_{1},\ldots,X_{n})$ be a homogeneous
polynomial of degree $d$ over $K$. Suppose $d<n$. Then there exists a point
$(x_{1},\ldots,x_{n})\in K^{n}$ distinct from the origin $(0,\ldots,0)$ such
that $F(x_{1},\ldots,x_{n})=0$.
\end{theorem}

Given a field $K$ and an integer $j$, one says that $K$ is a
$C_{j}$-field if every homogeneous polynomial over $K$ of degree $d$ in $n$
variables, such that $n>d^{j}$, has a non-trivial zero (i.e.\ besides the
origin) in $K^{n}$. The $C_{0}$-fields are the algebraically closed fields.
Chevalley's theorem asserts that finite fields are $C_{1}$ (one also calls
$C_{1}$ fields ``quasi-algebraically closed''). One can show that, if $K$ is
a $C_{j}$-field, the field $K(t)$ of rational functions in one variable over
$K$ and the field $K((t))$ of formal power series in one variable over $K$
are $C_{j+1}$-fields (\hyperref[bib:5]{[5]}). It had been, for a long time, an open question
whether $p$-adic fields are $C_{2}$. Recently this has been shown not to be
the case (\hyperref[bib:8]{[8]}).

\begin{proof}[Proof of Theorem \ref{thm:I.7.2}]
Let us write $q$ for $\operatorname{card}(K)$ and $p$ for the characteristic
of $K$ (so $q=p^{s}$). Let $V\subset K^{n}$ be the set of zeros of $F$,
i.e.\ the points $(x_{1},\ldots,x_{n})\in K^{n}$ such that $F(x)=0$ (we use,
hereafter, the symbol $x$ to stand for a point
$(x_{1},\ldots,x_{n})\in K^{n}$). According to Theorem \ref{thm:I.7.1}, (c),
$F(x)^{q-1}=0$ for $x\in V$ and $F(x)^{q-1}=1$ for $x\in K^{n}-V$. Thus the
polynomial $G(x)=F(x)^{q-1}$ is the characteristic function of $K^{n}-V$
with values in $\mathbf{F}_{p}$. The number modulo $p$ of points of
$K^{n}-V$ will thus be given by the sum $\sum_{x\in K^{n}}G(x)$. We are
going to calculate this sum and show that it is zero. It will follow that
$\operatorname{card}(K^{n}-V)$ is a multiple of $p$; inasmuch as
$\operatorname{card}(K^{n})=q^{n}=p^{ns}$ is also a multiple of $p$,
$\operatorname{card}(V)$ will also have to be a multiple of $p$. Certainly
$V$ contains the origin; so, if $p\mid\operatorname{card}(V)$, $V$
necessarily contains other points, since $p\geq2$. Thus, to prove Theorem
\ref{thm:I.7.2}, it suffices to show that
$\sum_{x\in K^{n}}G(x)=0\in\mathbf{F}_{p}$. Now, to calculate
$\sum_{x\in K^{n}}G(x)$, we observe that the polynomial $G$ is a linear
combination of monomials $M_{\alpha}(X)=X_{1}^{\alpha_{1}}\cdots
X_{n}^{\alpha_{n}}$. To determine $\sum_{x\in K^{n}}G(x)$ it suffices to
calculate
\[
\sum_{x\in K^{n}}M_{\alpha}(x)=\sum_{x\in K^{n}}
x_{1}^{\alpha_{1}}\cdots x_{n}^{\alpha_{n}}
=\left(\sum_{x_{1}\in K}x_{1}^{\alpha_{1}}\right)\cdots
\left(\sum_{x_{n}\in K}x_{n}^{\alpha_{n}}\right).
\]
The problem reduces to that of calculating sums of the form
$\sum_{z\in K}z^{\beta}$ ($\beta\in\mathbf{N}$).

(a) For $\beta=0$, $z^{\beta}=1$ for all $z\in K$. Consequently,
$\sum_{z\in K}z^{\beta}=\sum_{z\in K}1=q=0$.

(b) For $\beta>0$, the term $0^{\beta}$ is $0$, so the sum reduces to
$\sum_{z\in K^{*}}z^{\beta}$. $K^{*}$ is a cyclic group of order $q-1$
(Theorem \ref{thm:I.7.1}, (b)). Let $\omega$ generate $K^{*}$. Then
$\sum_{z\in K^{*}}z^{\beta}=\sum_{j=0}^{q-2}\omega^{\beta j}$ which is the
sum of a geometric progression. Thus:

(b') If $\omega^{\beta}\neq1$, i.e.\ if $\beta$ is not a multiple of $q-1$,
then
\[
\sum_{j=0}^{q-2}\omega^{\beta j}
=\frac{\omega^{\beta(q-1)}-1}{\omega^{\beta}-1}=0,
\]
since $\omega^{q-1}=1$.

(b'') If $\omega^{\beta}=1$, i.e.\ if $\beta$ is a multiple of $q-1$, then
\[
\sum_{j=0}^{q-2}\omega^{\beta j}=q-1.
\]

It follows from (a), (b'), and (b'') that
$\sum_{x\in K^{n}}x_{1}^{\alpha_{1}}\cdots x_{n}^{\alpha_{n}}$ vanishes
unless all the $\alpha_{i}$'s are non-zero and multiples of $q-1$. In this
case, the degree $\alpha_{1}+\cdots+\alpha_{n}$ of the monomial is at least
$(q-1)n$. However, since $G=F^{q-1}$, $G$ has degree $(q-1)d$ and
$(q-1)d<(q-1)n$ by assumption. Thus
$\sum_{x\in K^{n}}M_{\alpha}(x)=0$ for every monomial $M_{\alpha}(X)$ which
appears in $G$ with a non-zero coefficient. Therefore,
$\sum_{x\in K^{n}}G(x)=0$. We have seen that this relation implies the
theorem.
\end{proof}

Let us remark that it would have been sufficient, in place of the assumption
that $F$ was homogeneous, to have considered $F$ with no constant term.
Naturally, the strict inequality $d<n$ between the degree and the number of
variables is essential. For example, the norm of $\mathbf{F}_{q^{n}}$ to
$\mathbf{F}_{q}$ (cf.\ \hyperref[sec:2.6]{Chapter II, §6}) is a homogeneous polynomial of degree
$n$ in $n$ variables over $\mathbf{F}_{q}$ with no non-trivial zero.

\begin{statement}{Example}
A quadratic form in three variables over a finite field ``represents
zero'' (i.e.\ has a non-trivial zero). Passing from $K^{3}$ to the
projective plane $\mathbf{P}_{2}(K)$, this means that a conic over $K$
contains a point rational over $K$ (i.e.\ whose coordinates may be chosen in
$K$). The example of the conic $X^{2}+Y^{2}+Z^{2}=0$ over $\mathbf{R}$
(respectively $X^{2}+Y^{2}-3Z^{2}=0$ over $\mathbf{Q}$; in order to see that
$X^{2}+Y^{2}-3Z^{2}=0$ has no non-trivial solution in $\mathbf{Q}$, it
suffices to consider the case where $x,y,z$ are relatively prime integers
and then to reduce mod $4$) shows that we are not dealing with a property
common to all fields.
\end{statement}
